\documentclass[11pt]{article}
\usepackage{docstyle}
\renewcommand\sheetname{Classwork 1}

\begin{document}
\sheethead{Classwork 1}{Powers, factorising and algebraic fractions}{Grades: [A] [M] [E]}

Show all working. Leave answers with positive indices and fractions in simplest form.

\section*{Part A \textperiodcentered{} Powers, roots and factorising}
\qpart{A1}Simplify $\dfrac{(5x^{2})^{2}}{10x^{-3}}$, leaving your answer with positive indices.\lvl{A}
\af{\frac{25x^{4}}{10x^{-3}} = \frac{5x^{4-(-3)}}{2} = \frac{5x^{7}}{2}}
\al{}

\qpart{A2}Simplify $\sqrt[3]{125p^{9}}$.\lvl{A}
\al{$\sqrt[3]{125} = 5$,\quad $9 \div 3 = 3$,\quad so $5p^{3}$}

\qpart{A3}Factorise completely $9x^{2}y - 6xy^{2}$.\lvl{A}
\al{HCF $3xy$:\quad $3xy(3x - 2y)$}

\qpart{A4}Factorise $2x^{2} - x - 15$.\lvl{A}
\al{$ac = -30$; $-6$ and $5$:\quad $2x^{2} - 6x + 5x - 15 = 2x(x-3) + 5(x-3)$}
\al{$= (x-3)(2x+5)$}

\qpart{A5}Simplify $\dfrac{x^{2}-16}{x^{2}+2x-8}$.\lvl{A}
\af{\frac{(x-4)(x+4)}{(x+4)(x-2)} = \frac{x-4}{x-2}}

\section*{Part B \textperiodcentered{} Combine, solve, change the subject}
\qpart{B1}Write $\dfrac{3}{x+2} - \dfrac{x-4}{x^{2}+2x}$ as a single fraction in its simplest form.\lvl{M}
\af{\frac{3x - (x-4)}{x(x+2)} = \frac{2x+4}{x(x+2)} = \frac{2(x+2)}{x(x+2)} = \frac{2}{x}}
\al{u: common denominator $x(x+2)$ and correct numerator;\quad r: correct fraction, simplified fully}

\qpart{B2}Solve $\dfrac{3}{x} + \dfrac{1}{x-2} = 2$.\lvl{M}
\al{Multiply every term by $x(x-2)$:\quad $3(x-2) + x = 2x(x-2)$}
\al{$4x - 6 = 2x^{2} - 4x$,\quad so\quad $2x^{2} - 8x + 6 = 0$,\quad $x^{2} - 4x + 3 = 0$}
\al{$(x-1)(x-3) = 0$,\quad $x = 1$ or $x = 3$\qquad u: quadratic formed;\quad r: both solutions}

\qpart{B3}Make $t$ the subject of $m = \dfrac{3t-1}{t+2}$.\lvl{M}
\al{$mt + 2m = 3t - 1$,\quad $mt - 3t = -1 - 2m$,\quad $t(m - 3) = -(1 + 2m)$}
\af{t = \frac{1+2m}{3-m}\qquad \Bigl(\text{or } t = \frac{-1-2m}{m-3}\Bigr)}

\section*{Part C \textperiodcentered{} Show that}
\qpart{C}$x$ is a whole number. Show that $(x+1)^{3} - (x-1)^{3}$ is always \textbf{2 more than a multiple of 6}.\lvl{E}
\begin{plainbox}[title={Words you may use}]
expand \textperiodcentered{} collect like terms \textperiodcentered{} multiple of 6 \textperiodcentered{} whole number \textperiodcentered{} therefore
\end{plainbox}
Sentence starters (optional): \emph{Expanding the first bracket gives \ldots{}} \quad \emph{Subtracting, \ldots{}} \quad \emph{Since $x^{2}$ is a whole number, \ldots}
\al{$(x+1)^{3} = x^{3} + 3x^{2} + 3x + 1$ \hfill MP1}
\al{$(x-1)^{3} = x^{3} - 3x^{2} + 3x - 1$ \hfill MP1}
\al{$(x+1)^{3} - (x-1)^{3} = 6x^{2} + 2$ \hfill MP2}
\al{$x^{2}$ is a whole number, so $6x^{2}$ is a multiple of 6; therefore the result is 2 more \hfill MP3}
\al{than a multiple of 6, for every whole number $x$.}
\ansnote{6mm}{\small MP1 and MP2 together are r; MP3, a conclusion that names why $6x^{2}$ is a multiple of 6, makes it t.}

\section*{Extension}
Find every whole number $x$ for which $(x+1)^{3} - (x-1)^{3} = 56$.
\al{$6x^{2} + 2 = 56$,\quad $x^{2} = 9$,\quad $x = 3$ or $x = -3$}
\al{}
\end{document}
